\section{Prospective evaluation and verification}\label{sec:evidence}

\subsection{The earlier result is retained}

The first experiment completed 316 scheduled records. Its 24 selected
application models included 20 admitted models; baseline and matched
reformulation control solved 19, while both cut modes solved 18.
Automatic selection reduced direction LPs from 1,073 to 323 and callback
time from 5.59 to 3.05 seconds relative to unconditional activation,
but did not recover the lost solve. On 13 synthetic cases, the control
and cut modes solved all 13 while baseline solved 12, so the extra solve
could already be explained by reformulation.

All 1,082 recorded added cuts passed the first campaign's replay and all
270 returned incumbents passed its numerical residual checks. Eight
construction-error records lacked model/cut logs; those missing logs were
excluded from successful replay coverage and retained as unknown counts.
These historical results justify investigating a different integration,
not a favorable claim about its performance in advance.

\subsection{Frozen new comparison}

The new
\pathref{../experiments/protocol.md}{prospective protocol}
was fixed before primary outcomes. It excludes the prior convexification
campaign names and selects 30 cached MINLPLib OSiL models, ten each from
convex models, nonconvex continuous models, and nonconvex models with
integer variables. Eligibility uses size limits and syntax parsing, not
success of the new importer or cuts. Within each stratum, a fixed
SHA-256 ordering selects models. The 422 eligible models have stratum
counts 85, 208, and 129. ``New holdout'' refers to the named earlier
convexification campaigns, not to every use of these public instances
elsewhere in the repository.

Every selected failure remains in its denominator. The common builder
is used in \texttt{baseline}, \texttt{all}, and \texttt{auto}. Primary
full runs use seed zero, one SCIP thread, one BLAS/OpenMP thread, a
30-second soft total integration budget, and relative gap tolerance
$10^{-4}$. Model loading, analysis, building, cut generation, and solving
are included in the integration budget; process imports and independent
result checking are counted in outer wall time. The hard process limit
for these runs is 45 seconds. Nonpreemptive calls can overshoot a soft
budget and must be reported.

The fixed schedule comprises 90 new-holdout full runs, 30 historical
diagnostic runs, 39 full mechanism runs, 105 root-only runs, and 18
seed-one repeats: 282 jobs. The 13 mechanisms use ten-second soft
budgets; root-only runs use five seconds and one node. Mode order rotates
by model index. One worker runs at a time on a host shared with unrelated
jobs, so small timing differences remain descriptive. The 150-minute
aggregate emergency cap exceeds the sum of all scheduled hard process
timeouts; any unstarted jobs would still be recorded.

Frozen source copies and hashes, exact tagged parsed models, original
OSiL files, mode configurations, solver status, numerical bounds, nodes,
original-variable incumbents, work statistics, and cut witnesses are
retained. Every returned incumbent is evaluated against original rows,
bounds, domains, integrality, and objective by a separate evaluator at a
scaled tolerance of $10^{-5}$. Recorded root and final bounds are compared
with known synthetic optima or archived feasible MINLPLib objectives.
Such a conflict is a diagnostic flag, not by itself a proof about which
numerical calculation is wrong. Fresh-process replay checks each saved
cut against its original model and actual stored row, with tampering
controls and explicit accounting for missing logs.

The campaign source snapshot is distinct from the final working tree.
The first post-freeze changes rejected cut sides at SCIP's infinity
sentinel, checked that side in stored-row auditing, and caught failed
heuristic exchange-sample evaluations. The campaign then exposed four
worker errors in cut modes on \texttt{chp\_partload} and
\texttt{waterno2\_06}: affine splitting passed an unnecessarily large
list of symbolic generators to polynomial conversion. Discovery on a
large case also overran its allowance without checking between rows.
The repair uses each term's actual generators, safely declines terms
whose rational polynomial conversion is unavailable, and checks the
existing work allowance between discovery operations. These repairs have
their own targeted review and matched validation runs.

The archived prospective experiment keeps its source and every original
result. Its timings describe that frozen implementation. Missing worker
cut logs remain unknown. A separate supplement repeats the affected
cases with all comparison modes after the correction; it does not replace
the original outcomes or turn outcome-selected repairs into a new holdout.
The final reconciliation records exact repair selection, source identities,
work-budget effects, and replay coverage separately.

\subsection{Prospective campaign results}

The campaign completed all 282 scheduled jobs in 1,338.6 seconds of outer
wall time, with no hard process timeouts or unstarted jobs. There were
106 optimal statuses, 74 gap-limit statuses, 31 time limits, 67 node
limits, and the four retained worker errors described above. All 30 new
holdout models were admitted in every mode. The same 25 were numerically
solved in every mode; the same five remained unsolved. Here solved means
an optimal or gap-limit status with an incumbent passing the independent
original-model numerical checks.

\begin{center}
\begin{tabular}{lrrrr}
\toprule
Mode & Solved/selected & Added cuts & Cases with cuts & Total seconds\\
\midrule
Native baseline & 25/30 & 0 & 0 & 170.57\\
All admitted blocks & 25/30 & 38 & 7 & 187.02\\
Automatic selection & 25/30 & 10 & 3 & 187.12\\
\bottomrule
\end{tabular}
\end{center}
Total seconds sum the recorded integration and preparation time; they
exclude interpreter imports and independent result checking. Shared-host
timing is descriptive. The cut modes establish no solved-count gain in
this population and take more summed time in these runs.

Automatic selection reduced candidate LPs from 186 to 66 and support
calls from 138 to 51. It did not substantially reduce callback time:
27.63 seconds for \texttt{all}, versus 27.00 for \texttt{auto}.
Discovery accounted for 24.41 and 24.63 seconds respectively, preceding
most of the admission policy's savings. Of 30 full runs per cut mode,
28 invoked discovery, 12 found no supported blocks, and two solved
without calling this separator. The admitted groups contained 156
possibly overlapping blocks, of which 63 satisfied the structural
automatic rule; 186 source sides were declined. These are observed
coverage and work counts, not unique variables or a causal decomposition
of native relaxation strength.

At the prespecified scaled dual-comparison tolerance $10^{-6}$, final
bounds for \texttt{all} were better on two models, tied on 21, and worse
on seven; \texttt{auto} had two better, 24 tied, and four worse. Both
improved bounds occurred on models already solved in all modes.
Among the five common unsolved cases, both cut modes had zero better,
one tied, and four worse bounds. None of those five received an added
cut. For example, \texttt{graphpart\_clique-40} finished with lower
bound 438 for baseline and 392 for each cut mode. Discovery and callback
cost, time limits, and shared-host variation prevent attributing every
difference to one cause; these results do not show useful added strength
on the common hard cases.

The root-only application comparison gave one better, 26 tied, and three
worse bounds for \texttt{all}, versus one better, 28 tied, and one worse
for \texttt{auto}. The seed-one repeats solved all six
selected models in every mode and had six tied final bounds in each
comparison. The complete root and repeat tables remain in
\pathref{../experiments/campaign-v2/summary.json}{the campaign summary}.
The repetitions are too small to establish a population-wide timing effect.

Every mode solved 12 of the 13 full synthetic cases. A positive local
mechanism is retained: on the root-only \texttt{simplex\_quadratic\_vector}
case, both cut modes added two cuts and reached an optimal status with
bound approximately $-0.50000002$, versus baseline's node-limit bound
$-0.50027374$. Its true optimum is $-1/2$. This is evidence that the
row-domain support can supply useful information on a constructed case;
it is not an extra full-solve success on the application sample.

The ten historical diagnostics gave five solved cases in each mode.
The four cut-mode failures occurred after model construction, so they
must not be reported as importer refusals. The successful admission
checks for the seven earlier refused models likewise establish importer
coverage, not a guarantee that every subsequent separator operation
finishes successfully. The repaired runs below address those errors
and the discovery-budget defect directly.

Fresh-process replay passed all 123 recorded cuts, including their actual
stored SCIP rows, and verified 156 frozen source/input hashes. All 14
tampering controls were rejected. The four error records have no complete
cut logs and retain unknown counts; the successful replay claim excludes
them. All 271 returned incumbents passed the independent numerical
original-model checks, with no recorded reference-bound conflicts. The
\pathref{../experiments/campaign-v2/replay.json}{replay record} and
independent experiment audit provide the exact coverage. These results
certify the recorded cut inequalities under the documented trust base,
not the numerical solver's final global bounds or complete search.

\subsection{Matched repair validation}

The correction cohort was fixed before any repaired optimization outcome.
It includes every original model/phase/seed group whose cut mode exceeded
its recorded discovery allowance or failed on the identified global-symbol
recursion path. The rule selected 25 groups and 75 matched jobs, including
all three modes for each group. The exact triggering values and original
record hashes are retained in
\pathref{../experiments/repair-plan.json}{the repair plan}; the
\pathref{../experiments/repair-protocol.md}{repair protocol}
keeps the original phase, seed, time and node limits, and mode order.
The corrected code has a separate source snapshot. No activation
threshold, instance-name filter, or selected population is retuned.

This is a regression cohort selected from observed failures and budget
violations, not another prospective holdout. Its results remain separate
and cannot be spliced into the 30-model table above. All 75 jobs completed
in 754.5 seconds, with no worker errors, missing cut logs, hard timeouts,
or unstarted jobs. The solved counts were identical across the three
modes within each cohort:

\begin{center}
\begin{tabular}{lrrrr}
\toprule
Matched cohort & Groups & Solved per mode & All cuts & Auto cuts\\
\midrule
Full application runs & 8 & 4 & 12 & 3\\
Full historical diagnostics & 7 & 4 & 10 & 0\\
Root application runs & 9 & 0 & 11 & 3\\
Seed-one repeat & 1 & 1 & 3 & 0\\
\bottomrule
\end{tabular}
\end{center}
Every group has a baseline, an \texttt{all}, and an \texttt{auto} run;
the table's solved column applies to each mode, not their sum. Baseline
adds no new cuts. There are no additional solves from either cut mode.
For the eight affected full application models, final bounds were zero
better, five tied, and three worse for \texttt{all}; \texttt{auto}
had zero better, six tied, and two worse. Summed integration times were
123.91, 124.86, and 124.82 seconds for baseline, \texttt{all}, and
\texttt{auto}. For the seven historical diagnostics both cut modes had
six tied bounds and one worse bound, on \texttt{waterno2\_06}.
All root and seed-repeat bound comparisons tied at the stated tolerance.

The four formerly crashing diagnostic runs now stop incomplete discovery
and continue native SCIP. Six discovery calls in the correction cohort
crossed their soft deadline by a single nonpreemptive unit of work;
all six reported incomplete discovery. The largest discovery excess
was 3.618 milliseconds. The largest full-callback excess was 29.529
milliseconds, and the largest total soft-budget excess was 23.611
milliseconds. These measurements support the checked stopping contract,
not a hard real-time bound or a claim that the discarded large blocks
were successfully analyzed.

All 42 recorded repair cuts passed fresh-process original-model and actual-row
replay, with 103 archived source/input hashes verified and all 14 tampering
controls rejected. All 67 returned incumbents passed the independent numerical
checks, with no reference-bound conflicts. The
\href{run:../experiments/repair-discovery-v1/results.md}{repair results},
\href{run:../experiments/repair-discovery-v1/replay.json}{replay}, and
\href{run:../reviews/experiment-review.md}{independent experiment review}
retain the original and repaired evidence separately. Across the two
current campaigns, all 165 recorded cuts passed replay; the four original
unknown logs remain outside that claim.

The operational recommendation remains native SCIP by default. The
corrected implementation resolves the observed execution and discovery-budget
defects, supplies optional certified cuts, and has no demonstrated general
solve advantage in these experiments. The exact support and complete
tolerance-separation APIs remain reusable mathematical components independent
of that negative activation result.

\subsection{Targeted implementation and document checks}

The initial combined topic-specific local test command passed 250 tests and
eight subtests. After the discovery repair, the focused integration, row
arithmetic, and replay command passed 81 tests and three subtests. The
\href{run:../VERIFICATION.md}{verification record}
gives the actual command and preserves an earlier test-collection import
failure and its correction. Separate independent support diagnostics
checked 40 analytic and adversarial polytope cases, including equality
domains, flat valleys, affine images of simplexes, and exhaustive Boolean
MaxCut controls. The separation review checked 21 cases covering finite-net
completion, missed interior normals, thin domains, coordinate cones,
polynomial support error, tiny rational gaps, and altered certificates.
The numerical direction LP is not an oracle for any of these checks.

The source/domain review includes cancelled restrictions, one-sided
unbounded intervals, original affine implications, positive-base variable
powers, generic-expression constant handling, and deliberate model-binding
corruption. Integration replay tests reconstruct actual SCIP rows from
original source sides and reject altered premises. Owner tests and internal
independent reviews cover different risks; agreement of producer and replay
alone would not establish the underlying theorem.

The report's \href{run:VERIFICATION.md}{document verification record}
gives its targeted LaTeX build and evidence-link checks. No project-wide
verification or CI inspection is part of these local records. Finite tests
support implementation confidence under the proved contracts; they do not
prove a whole numerical SCIP solve or a general performance advantage.
